% ============================================================
% Section 6: Evaluation
% ============================================================
% Target length: ~1.5 columns (800–1000 words)
% Write: Month 6–7 (all team members)
% THIS IS THE MOST IMPORTANT SECTION FOR ACCEPTANCE
%
% DATA COLLECTION INSTRUCTIONS (see paper/data/ directory):
% All raw data goes in paper/data/ as CSV files.
% All measurement scripts go in Browser/scripts/bench/
% ============================================================

\section{Evaluation}
\label{sec:evaluation}

We evaluate \ferrite{} against four research questions:

\begin{enumerate}
  \item[\textbf{RQ1}] Does \ferrite{} enforce capability governance
    with 100\% coverage against realistic attacks?
  \item[\textbf{RQ2}] What is the performance overhead of the
    capability broker?
  \item[\textbf{RQ3}] Does the audit log maintain cryptographic
    integrity under adversarial conditions?
  \item[\textbf{RQ4}] Does content blocking provide measurable
    bandwidth savings?
\end{enumerate}

\subsection{Experimental Setup}

\todo{Describe the machine used for benchmarks.
      Required fields: CPU model, core count, RAM, OS,
      Rust version, Servo version.
      Also describe the test environment: network conditions,
      whether tests are run in Docker or on host.}

% DATA TO COLLECT FOR THIS SUBSECTION:
% - Run: systeminfo (Windows) or uname -a + lscpu (Linux)
% - Record: CPU, RAM, OS, kernel version
% - Save to: paper/data/system_info.txt

\subsection{RQ1: Enforcement Coverage}

% ── How to collect this data ────────────────────────────────
% Run: cargo test --workspace 2>&1 | tee ../paper/data/test_results.txt
% The adversarial test suite in:
%   crates/ferrite-shell/tests/broker_adversarial.rs
%   crates/ferrite-audit-log/tests/audit_adversarial.rs
% Each test = one adversarial scenario.
% Document each scenario in the table below.
% Add attack reproduction tests for the 4 real-world attacks.

\todo{Table: adversarial scenario results.
      Columns: Scenario ID | Attack Type | Component Tested |
               Expected Outcome | Result | Notes
      Target: 20+ rows.
      Real-world attack reproductions must appear as named rows:
      CometJacking, Screenshot Injection, Tainted Memories, HashJack.}

\begin{table}[h]
  \centering
  \caption{Adversarial scenario results. \cmark{} = blocked as expected.}
  \label{tab:adversarial}
  \begin{tabular}{llll}
    \toprule
    ID & Attack Scenario & Component & Result \\
    \midrule
    A1  & Expired token     & Broker & \todo{\cmark} \\
    A2  & Origin mismatch   & Broker & \todo{\cmark} \\
    A3  & Policy rejection (DomWrite) & Policy & \todo{\cmark} \\
    A4  & Token revocation  & Broker & \todo{\cmark} \\
    A5  & Unknown token ID  & Broker & \todo{\cmark} \\
    A6  & WebContent principal & Policy & \todo{\cmark} \\
    A7  & CometJacking repro & Broker+Policy & \todo{\cmark} \\
    A8  & Screenshot Injection repro & Broker & \todo{\cmark} \\
    A9  & Tainted Memories repro & Broker+Audit & \todo{\cmark} \\
    A10 & HashJack repro    & Broker & \todo{\cmark} \\
    \todo{...} & \todo{10+ more} & & \\
    \bottomrule
  \end{tabular}
\end{table}

\subsection{RQ2: Broker Performance}

% ── How to collect this data ────────────────────────────────
% Add a benchmark binary: crates/ferrite-shell/src/bin/bench_broker.rs
% Measure:
%   - 10,000 iterations of broker.check() with a valid token
%   - 10,000 iterations with an expired token (fast path)
%   - 10,000 iterations with policy evaluation (regorus)
% Record: min, median, p95, p99, max in microseconds
% Also measure with 1, 10, 50 active tokens in the HashMap
% Save raw output to: paper/data/broker_latency.csv
% Format: scenario,iteration,latency_us
%
% Run on BOTH debug and release builds.
% Report release build numbers in the paper.
%
% Command:
%   cargo build --release -p ferrite-shell
%   cargo run --release -p ferrite-shell bench 2>&1 | tee ../paper/data/broker_latency_raw.txt

\todo{Figure: CDF or box plot of broker latency distribution.
      X-axis: latency in microseconds.
      Y-axis: cumulative fraction of requests.
      Show three lines: no-policy, allow-all policy, deny-by-default policy.
      Generate with: python3 scripts/plot_latency.py paper/data/broker_latency.csv}

\begin{table}[h]
  \centering
  \caption{Capability broker latency (\textmu s), release build.
           \todo{Fill from paper/data/broker\_latency.csv}}
  \label{tab:latency}
  \begin{tabular}{lrrrrr}
    \toprule
    Scenario & Min & Median & P95 & P99 & Max \\
    \midrule
    No policy     & \todo{} & \todo{} & \todo{} & \todo{} & \todo{} \\
    Allow-all     & \todo{} & \todo{} & \todo{} & \todo{} & \todo{} \\
    Deny-by-default & \todo{} & \todo{} & \todo{} & \todo{} & \todo{} \\
    \bottomrule
  \end{tabular}
\end{table}

\subsection{RQ3: Audit Log Integrity}

% ── How to collect this data ────────────────────────────────
% Tests in: crates/ferrite-audit-log/tests/audit_adversarial.rs
% Specifically:
%   - test_tampered_entry_hash_detected
%   - test_tampered_prev_hash_detected
%   - test_sqlite_detects_tampered_chain_on_load
%   - test_100_entries_chain_valid
% Record: verify_chain() result, time to verify N entries
% Save to: paper/data/audit_integrity.csv
%
% Also measure: time to verify chains of 100, 1000, 10000 entries
% This shows the log scales correctly.

\todo{Table or figure: chain verification time vs number of entries.
      Show it scales linearly (expected for hash chain).
      Show that any single-entry tamper is detected.}

\subsection{RQ4: Bandwidth Savings}

% ── How to collect this data ────────────────────────────────
% Test URLs (tracker-heavy sites):
%   cnn.com, bbc.com, reddit.com, theguardian.com, forbes.com,
%   huffpost.com, dailymail.co.uk, usatoday.com, nypost.com, msn.com
%
% Method:
%   1. Disable adblock-rust in Servo (feature flag or config)
%   2. Load each URL, record total bytes from Servo network stats
%   3. Enable adblock-rust
%   4. Load each URL again, record total bytes
%   5. Compute % reduction per site
%   Save to: paper/data/bandwidth_savings.csv
%   Format: url,bytes_without_blocking,bytes_with_blocking,reduction_pct
%
% Note: control for caching by clearing between runs
% Run each URL 3 times, take median

\todo{Figure: bar chart of bandwidth savings per site.
      Show \% reduction on Y-axis, site name on X-axis.
      Mark the 10\% target line.}

\subsection{Discussion}

\todo{3--4 paragraphs:
      1. What the results mean for the research questions
      2. Honest discussion of limitations (Servo compat, etc.)
      3. Comparison to model-level defenses (cite Nasr et al.)
      4. The key claim: architectural enforcement works
         independently of model behavior}
